% v20. Complete G5_11 obstruction.  No denominator or scale cutoff.
\section{\texorpdfstring{$G_{5,11}$}{G5,11} 型の無上界排除}
\label{g51120:section}

\begin{theorem}[二つの三価頂点を持つ七高辺型の障害]
\label{g51120:no-support}
五次の算術的接着データの高分母 support は
\[
 H=G_{5,11}=\{01,02,03,04,12,13,24\}
\]
と同型ではあり得ない．分母の上限は仮定しない．
正の整数局所 charts，marked clutching，非負 node，対向 $D$ の相殺を用いる．
\end{theorem}

この graph では $3,4$ が二価，$1,2$ が三価，$0$ が四価である．
低辺は $14,34,23$ であり，$0$ は二つの二価頂点の共通高近傍である．
証明では最初に二価 cores の reverse data を数学的に有限化する．
その後，分母と scale の全無限族を保持した厳密な整数演算で排除する．

\subsection{四つの向きと単位 entry の禁止}

二価 core の canonical order $(A,B,C,D)$ では
\begin{equation}\label{g51120:canonical}
 P=\begin{pmatrix}
 0&-U&x&-s\\
 -V&0&-y&z\\
 r&-r-1&0&0\\
 -S&T&-L&0
 \end{pmatrix},
\end{equation}
\[
 \begin{gathered}
 a=sS+1,\quad b=zT+1,\quad x=sS/r,\quad y=zT/(r+1),
 \\
 U=a+x-s,\quad V=b-y+z,\quad L=T-S+1.
 \end{gathered}
\]
全ての $r,s,z,S,T,x,y$ は正整数で，$a,b$ は奇数，$UV=ab-1$ である．
既知の二分補題により，正の枝は
\begin{equation}\label{g51120:positive}
 S>r,\qquad rT>(r+1)S,\qquad L\ge3,
\end{equation}
負の枝は
\begin{equation}\label{g51120:negative}
 S=T=w\ge2,\quad r=w+H,\quad
 \frac{w^2}{2}<H(H+1)<2w^2
\end{equation}
を満たす．特に $S\ge2$ である．正の枝では $T\ge5$
（補題~\ref{cl16:T-five}）である．

\begin{lemma}[正の単位 entry の一つの禁止]\label{g51120:no-x-one}
式~\eqref{g51120:canonical} では $x=1$ は起こらない．
\end{lemma}
\begin{proof}
正の枝では $x=sS/r>1$ である．負の枝で $x=1$ と仮定すると
$r=sS$，$s\ge2$，$T=S$ である．高高 pair equation を代入して整理すると
\[
 z\{(-s^2+3s-1)S+2-s\}+(2-s)(sS+1)=0
\]
を得る．$s=2$ では左辺は $zS>0$，$s\ge3$ では両項が負であり，
どちらも不可能である．
\end{proof}

二価 core を $i\in\{3,4\}$，もう一方を $j$，
$i$ 固有の高近傍を $h$，残る三価頂点を $k$ と書く．
すなわち $(i,j,h,k)=(3,4,1,2)$ または $(4,3,2,1)$ である．
canonical 標識の四通りを次表で固定する．最後の欄は ray $j$ の行を
列 $(0,h,k)$ の順で記したものである．
\begin{center}\small
\begin{tabular}{ccccccl}
型&$A$&$B$&$C$&$D$& &$(p^i_{j0},p^i_{jh},p^i_{jk})$\\\hline
I&$0$&$h$&$j$&$k$&&$(r,-r-1,0)$\\
II&$h$&$0$&$j$&$k$&&$(-r-1,r,0)$\\
III&$0$&$h$&$k$&$j$&&$(-S,T,-L)$\\
IV&$h$&$0$&$k$&$j$&&$(T,-S,-L)$
\end{tabular}
\end{center}

共有高辺 $ih$ の共通低 puncture $j$ で
\[
 u=p^i_{jh},\qquad v=p^h_{ji}
\]
と置く．補題~\ref{ls11:low-puncture} を reverse entries に書き換えると
\begin{equation}\label{g51120:partner}
 v=-u,\qquad \{u,v\}=\{1,q_{ih}-1\},\qquad
 \{u,v\}=\{-1,-(q_{ih}-1)\}
\end{equation}
の三択である．通常型 $v=-u$ では $|u|\le4$
（補題~\ref{cl16:ordinary-bound}）である．また $u<0$ なら，
相手の二価 core $j$ における $p^j_{ih}$ が低低 entry で非正なので
\[
 1-u-p^j_{ih}=n_{ij;h}\le5
\]
から例外型も含めて $|u|\le4$ となる．

\subsection{正の枝の III 型を先に排除する}

\begin{lemma}\label{g51120:no-positive-III}
二つの二価 cores のいずれも，正の枝の III 型には属さない．
\end{lemma}
\begin{proof}
対称性により core $3$ が正の III 型であるとする．このとき $u=T\ge5$ である．
低辺 $34$ の puncture $1$ の node 非負性から
$p^4_{31}\le1-T<0$ であり，core $4$ では $D=3,C=1$ となる．
従って $p^4_{13}=0$ であり，低辺 $14$ の node を用いると
\[
 v=p^1_{43}=1-n_{14;3}\le1.
\]
通常型は $T>4$ に反するので，式~\eqref{g51120:partner} は
$v=1$，$q_{13}=T+1$ を強制する．従って $z=1$，$n_{14;3}=0$ である．

三高 matching 定理~\ref{mh18:matching} より，core $1$ の低 row の残り二成分は負である．
core $4$ の $C=1$ の行は，高 rays に $r',-r'-1$ を持つ．
後者に対応する node は少なくとも $r'+3$ なので，総数五から $r'\le2$ を得る．
$S'\ge2$ より core $4$ は正の枝である．

もし core $4$ が IV 型なら，低辺 $34$ の puncture $2$ の node は
$1+L+S'\ge6$ となり不可能である．従って core $4$ も III 型である．
puncture $0$ の node は $1+S+S'\le5$ なので $S=S'=2$，$r=r'=1$ となる．
残る二つの node は $T'-T,T-T'$ であり，$T=T'$ である．
同じ議論を反対側に適用すると $z'=1$ である．

$r=1,S=2,z=1$ の pair equation は
\[
 s(8-T)=T-4
\]
である．$T>4$ と $y=T/2\in\ZZ$ より $T$ は六以上の偶数なので，
$T=6,s=1$ に限られる．両 core は
\[
 (r,s,z,S,T)=(1,1,1,2,6),\qquad (a,b)=(3,7)
\]
となる．core $4$ の ray $1=C$ の係数は $D^4_1=5$ である．
core $1$ の低 row は，matching と node 総数五から
\[
 (p^1_{40},p^1_{42},p^1_{43})=(-1,-1,1)
\]
に強制される．pair equations による低 column の和から
$D^1_4=q_{01}+q_{12}-9$ を得るので，辺 $14$ の $D$ 相殺は
$q_{01}+q_{12}=4$ を要求する．両分母は三以上なので矛盾する．
\end{proof}

\subsection{全 reverse data に対する証明された有限上界}

\begin{proposition}\label{g51120:reverse-bounds}
残る二価 cores の canonical reverse data は全て
\begin{equation}\label{g51120:box}
 1\le r\le9,\qquad 2\le S\le7,\qquad S\le T\le13
\end{equation}
を満たす．さらに I 型では $r\le3$，II 型では $r\le4$，
IV 型では $S\le4$ であり，III 型は負の枝に限られる．
この上界は分母 $a,b$ または scales $s,z$ の上限を意味しない．
\end{proposition}

\begin{proof}
I 型は $u=-r-1<0$ なので $r\le3$ である．
補題~\ref{tn14:small-negative-branch} により，この範囲に負の枝はない．
正の枝では $y=zT/(r+1)>1$ なので，$u$ の負例外型は起こらず，
必ず $v=r+1$ の通常型である．

II 型で $r>4$ なら通常型は不可能であり，正例外型は
$v=1,q_{ih}=r+1$ を要求する．これは $x=1$ に反する．従って $r\le4$ である．
III 型は負の枝だけが残る．その $u=T=w$ が四より大きい場合も
正例外型が $q_{ih}=T+1$，すなわち $z=1$ を要求するが，
$y=T/(r+1)<1$ となって不可能である．従って $w\le4$ である．
IV 型では $u=-S$ なので $S\le4$ である．
以上から，どの型でも負の枝なら $w\le4$ であり，
式~\eqref{g51120:negative} は $H\le5$，$r=w+H\le9$ を与える．
以下では正の枝だけを調べる．

\paragraph{IV 型．}
相手が I または IV 型なら，両者の ray $34$ の行が高 ray $0$ に正成分を持ち，
$n_{34;0}<0$ となる．相手が II 型なら
$n_{34;0}=r'+2-T\ge0$ なので $T\le6$ である．
相手が III 型なら，既に負の枝で $S'=T'\le4$ であるから，
$n_{34;0}=1+S'-T\ge0$ により $T\le5$ である．
従って IV 型は $r<S\le4,T\le6$ を満たす．

\paragraph{I 型の相手と上界．}
相手が I または IV 型なら，やはり puncture $0$ の node が負になる．
相手が III 型なら，I 型の低低 entry が零であるため，
相手固有の高 puncture の node が $1-T'<0$ となる．
従って相手は II 型であり，その同じ node が $1-r'\ge0$ を与えるので $r'=1$ である．

core $3$ を I 型，core $4$ を II 型とする．core $1$ の低 row の成分
$v=p^1_{43}=r+1$ は正なので，残る二成分は負である．
core $4$ の ray $1=D$ における逆向き成分は $T',-S',-L'$ であり，
必要な node は
\[
 n_{14;3}=L'-r,\qquad n_{14;2}\ge S'+2.
\]
従って $S'+L'\le r+3$，すなわち $T'\le r+2\le5$ である．
正の枝の $T'\ge5$ と $r'=1$ を合わせると
\[
 r=3,\qquad (r',S',T')=(1,2,5)
\]
を得る．この相手の pair equation と整数性は
\[
 (s'-3)(z'+2)=-8,\qquad z'\in2\ZZ_{>0}
\]
であり，$(s',z')=(1,2),(2,6)$ に限られる．従って $q_{24}=3$ または $5$ である．

辺 $24$ の共通低 puncture $3$ では，相手の reverse が $r'=1$ なので，
$R=p^2_{34}$ は $-1$ または $q_{24}-1$ である．後者なら matching により
$p^2_{30}<0$ なので，辺 $23$ の node を $m$ と書くと
\[
 m_0\ge S+2,\qquad m_4=L+2-q_{24}.
\]
その総数五から $T\le q_{24}\le5$ を得るが，$r=3,S>3$ では $T\ge6$ であり矛盾する．
従って $R=-1$，$m_4=L+2\le5$ である．正の枝の $L\ge3$ より
$L=3$，$3<S<6$ となる．従って $(S,T)=(4,6),(5,7)$ だけが残る．

\paragraph{II 型と II 型．}
両者の共通高 ray $0$ への成分は $-r-1,-r'-1$ なので
$n_{34;0}=r+r'+3\le5$，従って $r=r'=1$ である．両者は正の枝に属する．
辺 $13$ の通常型は $p^1_{43}=-1$ を要求するが，
このとき $n_{14;3}=L'+2\le5$ なので $L'=3$ となる．
これは $r'=1$，$S'>1$，$r'(T'-S')>S'$ と両立しない．
従って両方の共通 punctures は正例外型である．matching と node 総数から
\[
 T'\le q_{13}\le T'-S'+3,\qquad
 T\le q_{24}\le T-S+3
\]
となり，$S,S'\le3$ が従う．
$r=1$ の pair equation は
\[
 z(s+1)T=(4S-2)sz+2(S-1)s+2z+2
\]
であり，$z\ge1$ から
\[
 T\le\frac{(6S-4)s+4}{s+1}<6S-4\le14
\]
を得る．従って $T\le13$ である．

\paragraph{II 型と IV 型．}
II 型の相手が III 型なら，その固有高 puncture の node は $1-T'<0$ である．
I 型との組は既に扱ったので，残るのは IV 型との組だけである．
core $3$ を正の II 型，core $4$ を IV 型とする．$S'\le4$ は既知である．
辺 $24$ の共通低 puncture $3$ では，core $4$ の reverse は $-S'<-1$ なので
$R=p^2_{34}$ は $S'$ または $-1$ に限られる．
通常の $R=S'>0$ なら，matching は $p^2_{31}<0$ を要求し，
\[
 m_1\ge S+2,\qquad m_4=L+1-S'
\]
から $T\le S'+1\le5$ を得る．負例外の $R=-1$ なら
$m_4=L+2\le5$ であり，正の枝より $L=3$ となる．
従って $r<S<2r\le8$，$S\le7,T\le9$ である．
これで四つの向きの全てと全ての相手の組を尽くし，主張が従う．
\end{proof}

\subsection{全 scale の点・半直線表示}

$h=S-r$，$\delta=T-S$ と置くと，高高 pair equation は
\begin{equation}\label{g51120:bilinear}
 Fsz+As+Bz+C=0
\end{equation}
に等しい．ただし
\[
 F=(2r+1)h-r^2(\delta-1),\quad A=h(r+1),\quad
 B=-r(h+\delta-1),\quad C=r(r+1).
\]
本節では $L\le4$ を仮定できないため，$F=0$ の場合も保持する．
$F\ne0$ なら
\begin{equation}\label{g51120:factor}
 (Fs+B)(Fz+A)=K,\qquad K=AB-FC
\end{equation}
である．$K\ne0$ ならその全 signed divisors を列挙すればよい．
$K=0$ なら $Fs+B=0$ または $Fz+A=0$ の全直線を残す．
$F=0$ では式~\eqref{g51120:bilinear} 自体が整数一次式になる．
負の枝の $F=0$ は $2r+1\mid r^2$ を要求するので起こらず，
この場合は $A>0,B<0$ であり，正の整数解を全等差半直線として表示できる．

各直線上で
\[
 r\mid sS,\qquad r+1\mid zT,\qquad
 sS\equiv zT\equiv0\pmod2
\]
を課す．例えば parameter を $2r(r+1)$ を法として分ければ全整数解を尽くす．
正値性で始点をずらすと，各族は
\begin{equation}\label{g51120:halfline}
 (s,z)=(s_0,z_0)+t(d_s,d_z),\qquad t\in\ZZ_{\ge0}
\end{equation}
となる．どの半直線も途中で打ち切らない．

命題~\ref{g51120:reverse-bounds} の範囲で，全ての点と半直線を作ると，
一つの二価 core あたり次の表になる．
\begin{center}\small
\begin{tabular}{crrr}
向き&点型&全半直線型&合計\\\hline
I&27&30&57\\
II&31&30&61\\
III&8&0&8\\
IV&24&18&42\\\hline
合計&90&78&168
\end{tabular}
\end{center}
これは整除・偶奇の剰余類による分割も含む表示数であり，幾何的模型の個数ではない．
無限族を持つ reverse triples は，次の四つだけである：
\begin{center}\small
\begin{tabular}{ccccl}
$r$&$S$&$T$&場合&全正整数 scale の表示\\\hline
1&2&6&$F=0$&$s=2z-1$, $z\ge1$\\
2&6&12&$F=0$&$s=3a-2$, $z=2a-1$, $a\ge1$\\
2&3&6&$K=0$&$s=2a$, $z=1$, $a\ge1$\\
3&4&6&$K=0$&$s=3a$, $z=2$, $a\ge1$
\end{tabular}
\end{center}
後二つは固定分母七・十三の共鳴族である．前二つの affine 族も残すことが必要である．

\subsection{低行の整数接着と最後の parity 障害}

各二価 core の族を一つずつ選び，まず低辺 $34$ の三つの node
\[
 n_{34;k}=1-p^3_{4k}-p^4_{3k},\qquad k=0,1,2
\]
が非負であるものを残す．これらは reverse data だけで定まり，
scale の自由 parameter には依存しない．この段階の型対は $3592$ 個である．

低辺 $14,23$ の node 分配をそれぞれ
\[
 (n_0,n_2,n_3),\qquad(m_0,m_1,m_4)
\]
とし，非負整数で各和が五である全二十一通りを使う．
三高 cores $1,2$ の唯一の低 row は
\begin{equation}\label{g51120:low-rows}
 R^1_k=1-n_k-p^4_{1k}\quad(k=0,2,3),\qquad
 R^2_k=1-m_k-p^3_{2k}\quad(k=0,1,4)
\end{equation}
である．全成分が非零で，正成分がちょうど一つであるものだけを残す．
さらに $(p^3_{41},R^1_3)$ と $(p^4_{32},R^2_4)$ に
式~\eqref{g51120:partner} の三択を課すと，$1395$ 配置が残る．

通常型は分母に新しい式を課さない．例外型では，二つの reverse entries の和の絶対値が
共有高分母 $q_{13}$ または $q_{24}$ に等しい．各分母は式~\eqref{g51120:halfline} の
parameter に関して一次式であり，この条件は正確な整数の等式として処理する．
定数の不一致で $759$ 配置，parameter が整数にならないことで $550$ 配置が消え，
合計 $1309$ 配置が排除される．残る $86$ 配置では，それぞれの parameter は
一つに固定されるか，全ての非負整数を動くままである．

最後に，低辺 $34$ の対向係数を足す．canonical な低 rays の係数は
\[
 D_C=T-S+y-x,\qquad D_D=s-z-1
\]
なので，その残差 $D^3_4+D^4_3$ は残る自由 parameters に関する一次式である．
全配置の結果は次表のとおりであり，零になるものはない．
\begin{center}\small
\begin{tabular}{crl}
向きの組&配置数&$D^3_4+D^4_3$ の障害\\\hline
I--II または II--I&32&全自由 parameters に対して奇数\\
II--II&5&四配置で $8$，一配置で $28$\\
II--IV または IV--II&40&全自由 parameters に対して奇数\\
III--III&9&常に $-4$\\\hline
合計&86&全配置で相殺不可能
\end{tabular}
\end{center}

最後の表は，node 配分と scale の剰余類表示の重複を除くと，
$3\leftrightarrow4$ の交換を除いて十組の canonical models へ圧縮できる．
五つ組を $(r,s,z,S,T)$ の順とし，
\[
 \begin{gathered}
 F_a=(3,3a,2,4,6)\quad(a\ge1),\\
 A=(1,1,2,2,5),\quad B=(1,2,6,2,5),\quad C=(1,3,2,2,6),\\
 E=(2,2,6,3,5),\quad N_0=(4,4,5,2,2),\\
 N_1=(7,7,2,4,4),\quad N_2=(7,21,4,4,4),\\
 N_3=(8,6,9,4,4),\quad N_4=(8,8,27,4,4)
 \end{gathered}
\]
と書くと，残差は次のようになる．
\begin{center}\small
\begin{tabular}{ccl}
一方の向き・模型&他方の向き・模型&$D^3_4+D^4_3$\\\hline
I, $F_a$&II, $A$&$11-4a$\\
I, $F_a$&II, $B$&$19-4a$\\
II, $F_a$&IV, $A$&$3-4a$\\
II, $B$&II, $B$&$28$\\
II, $C$&II, $C$&$8$\\
II, $E$&IV, $N_1$&$13$\\
II, $E$&IV, $N_2$&$25$\\
II, $E$&IV, $N_3$&$5$\\
II, $E$&IV, $N_4$&$-11$\\
III, $N_0$&III, $N_0$&$-4$
\end{tabular}
\end{center}
従って，残る唯一の無限族 $F_a$ についても，全正整数 $a$ に対して矛盾が成立する．

この有限表は，証明から得た reverse 上界と signed-divisor の完全因数分解によるものである．
分母や半直線 parameter の数値上限は使用していない．
\nolinkurl{verify_g5_11_obstruction_v20.py} は，全 $168$ 型の affine 行列と
全 $1395$ 配置を生成し，元の行和と pair products を parameter の多項式恒等式として確認する．
結果 JSON には，各 node 行，二つの低 rows，例外型の要求分母，parameter の全範囲，
および最後の奇数残差または非零定数を保存する．
別実装は主生成器を import せず，canonical データからこれらを再構成できる．

補題~\ref{g51120:no-positive-III} と最後の完全表により，四つの向きの全てが排除される．
従って定理~\ref{g51120:no-support} が従う．
